\documentclass[10pt,a4paper]{amsart}
\usepackage[margin=19mm]{geometry}
\usepackage{amsmath,amssymb,mathtools}
\usepackage[scr=rsfs,cal=euler]{mathalfa}
\usepackage{xfrac}
\usepackage[unicode,hidelinks,bookmarks=false]{hyperref}
\newcommand{\ie}{i.e.\ }
\newcommand{\cf}{cf.\ }
\newcommand{\resp}{respectively\ }
\newcommand{\Subsection}{\subsection}
\providecommand{\nobreakdash}{\nobreak\hbox{-}}
\setlength{\emergencystretch}{2em}
\multlinegap0pt
\usepackage{xcolor}
\usepackage{cleveref}
\theoremstyle{plain}
\newtheorem{theorem}{Theorem}
\newtheorem{proposition}{Proposition}
\DeclareMathOperator{\Var}{Var}
\DeclarePairedDelimiter{\abs}{\lvert}{\rvert}
\providecommand{\E}{\mathbb{E}}
\providecommand{\limsup}{\operatornamewithlimits{limsup}}
\providecommand{\liminf}{\operatornamewithlimits{liminf}}
\title[The martingale strong-law convergence rate proposition]{Convergence rate of the (1+1)-evolution strategy on locally strongly convex functions with Lipschitz continuous gradient: the martingale strong-law convergence rate proposition (Proposition 2)}
\author{Daiki Morinaga, Kazuto Fukuchi, Jun Sakuma, Youhei Akimoto}
\date{Source: arXiv:2209.12467v4 (CC BY 4.0)}
\begin{document}
\maketitle
\begin{abstract}
Evolution strategy (ES) is one of the promising classes of algorithms for black-box continuous optimization.
Despite its broad successes in applications, 
theoretical analysis on the speed of its convergence is limited on convex quadratic functions and their monotonic transformation.%theoretically how fast it converges to a optima on convex functions is still vague.
In this study, an upper bound and a lower bound of the rate of linear convergence of the (1+1)-ES on locally $L$-strongly convex functions with $U$-Lipschitz continuous gradient are derived as $\exp\left(-\Omega_{d\to\infty}\left(\frac{L}{d\cdot U}\right)\right)$ and $\exp\left(-\frac1d\right)$, respectively.
Notably, any prior knowledge on the mathematical properties of the objective function, such as Lipschitz constant, is not given to the algorithm, whereas the existing analyses of derivative-free optimization algorithms require it.
\end{abstract}

\noindent\textbf{Erratum:}
The version of this paper published at IEEE TEVC (2024) contains some technical errors. The errors do not affect the correctness
of the theorems. They are corrected in this version. For clarity, the changes are marked in {\color{blue}
blue}.


\noindent\emph{Source and licence.} Convergence rate of the (1+1)-evolution strategy on locally strongly convex functions with lipschitz continuous gradient. Version arXiv:2209.12467v4 (CC BY 4.0), distributed under the Creative Commons Attribution 4.0 International licence, http://creativecommons.org/licenses/by/4.0/, as recorded in the arXiv licence metadata for this exact version and shown on the abstract page https://arxiv.org/abs/2209.12467v4. Full licence notice: licenses/arxiv-2209.12467-CC-BY-4.0.txt.\par\smallskip

\noindent\textit{Editorial scope.} Counted unit: the article's convergence-rate proposition for Markov chains driven by a martingale difference sequence, printed as Proposition 2 (source0.tex lines 455--475, source label \texttt{prop:convergence\_rate}), delivered together with the article's complete original proof of it (source0.tex lines 2410--2425, the appendix subsection ``Proof of Proposition 2'' with its 16 lines of proof, with no gap and no deferral). No mathematics is written, repaired or re-derived: statement and proof are the article's own text line for line, in the article's own notation ($\{\mathcal{F}_t\}$, $\{X_t\}$, the conditions C1, C2, C3, $B^\mathrm{upper}$, $B^\mathrm{lower}$, $Z_t$, statements (I) and (II)), and no line of the retained spans is omitted, substituted or re-flowed; the commented-out display attempts inside the proposition (lines 465--468 and 471--474) are carried verbatim. Required context retained uncounted: the article's abstract (lines 186--192, reproduced verbatim as the wrapper's abstract); the article's Erratum notice (lines 178--181), moved to follow the abstract in the wrapper because the AMS class requires the abstract immediately after the title, and recording that the corrections relative to the published IEEE TEVC version are marked in blue, which is the key to the blue-marked terms in the retained statement and proof; the article's informal definition of linear convergence with rate $r$ (lines 426--429); and the article's sentences introducing the counted proposition and its appendix proof (lines 451--454), immediately followed by the article's appendix section heading ``Proof'' (line 2400) and the appendix subsection heading (line 2410) that the retained sentence cross-refers to. The retained text contains one blue-marked correction in condition C3 of the statement (line 461) and two in the proof (lines 2422--2423), retained literally with their \texttt{\textbackslash color\{blue\}} commands. Editorial formatting only, in addition: an AMS \texttt{amsart} preamble that restates the macros and environments the retained lines use (the operators \texttt{\textbackslash Var} and \texttt{\textbackslash E} as the article's own preamble declares them, the \texttt{\textbackslash abs} paired delimiter that the article declares with \texttt{mathtools}, the \texttt{\textbackslash limsup} and \texttt{\textbackslash liminf} declarations, the \texttt{proposition} environment as the article declares it (the wrapper gives it its own counter rather than the article's shared counter, so that the \texttt{cleveref} package names the retained cross-reference ``Proposition'' as the article does), the \texttt{xcolor} package for the blue correction marks and the \texttt{cleveref} package for the \texttt{\textbackslash Cref} commands of the retained headings); and one counter setting so that the retained proposition prints as Proposition 2, its number in the article. The article is typeset in the IEEEtran journal class with Times text and its own colours and the wrapper is an amsart document with Latin Modern fonts and black hyperlinks, so the visual presentation differs although the mathematical text does not. Bibliography: the retained text cites exactly one entry of the article's own in-source \texttt{thebibliography}, transcribed verbatim with the article's own printed number as its label ([7] chow1967strong); no other entry of the article's bibliography is reproduced, and no citation outside that one occurs in any retained line. No figure is retained: the article's 22 graphics inclusions (author photographs and numerical plots in its section on numerical evaluation) lie outside every retained span, the retained lines contain no \texttt{\textbackslash includegraphics} command, and the wrapper declares no asset.
A deterministic real-valued sequence $\{X_t\}_{t\geq 0}$ is said to convergence linearly with convergence rate $r \in (0, 1)$ if
$\lim_{t\to\infty} \abs*{X_{t+1} - X^\mathrm{opt}}/\abs*{X_t - X^\mathrm{opt}} = r$. A slightly more applicable definition is $\lim_{t\to\infty} \frac{1}{t}\log (\abs*{X_{t} - X^\mathrm{opt}} /\abs*{X_0 - X^\mathrm{opt}}) = \log(r)$.
For a stochastic real-valued sequence, various generalizations are considered. 
In this study, the definition of the convergence rate is the limit of the slope of the logarithmic distance between the current solution and the optimum averaged over time.

Our main mathematical tool to derive bounds for the upper and lower convergence rate of the (1+1)-ES is the strong law of large numbers on martingales \cite{chow1967strong}.
It reduces the analysis of the limit behavior to the analysis of the expected sequence change in one step. 
The following proposition is the central tool for our analysis.
Its proof is provided in \Cref{apdx:prop:convergence_rate}.

\setcounter{proposition}{1}
\begin{proposition}\label{prop:convergence_rate}
Let $\{\mathcal{F}_t\}_{t \geq 0}$ be a filtration of a $\sigma$-algebra, and $\{X_t\}_{t \geq 0}$ be a Markov chain adapted to $\{\mathcal{F}_t\}_{t \geq 0}$.
Consider the following conditions:
\begin{enumerate}
\item[C1] $\exists B^\mathrm{upper} > 0$ such that $\E[X_{t+1} - X_{t} \mid \mathcal{F}_t] \leq - B^\mathrm{upper}$ for all $t \geq 0$;
\item[C2] $\exists B^\mathrm{lower} > 0$ such that $\E[X_{t+1} - X_{t} \mid \mathcal{F}_t] \geq - B^\mathrm{lower}$ for all $t \geq 0$;
\item[C3] $\sum_{t=1}^{\infty}{\color{blue}\Var[X_{t}]} / t^2 < \infty$.
\end{enumerate}
(I) If C1 and C3 hold, then we have
$\Pr\left[\limsup_{t \to \infty} \frac{1}{t} (X_t - X_0) \leq - B^\mathrm{upper} \right] = 1$.
% \begin{align}
% \Pr\left[\limsup_{t \to \infty} \frac{1}{t} (X_t - X_0) \leq - B^\mathrm{upper} \right] = 1 \label{eq:prop:upper}
% .
% \end{align}
(II) If C2 and C3 hold, then we have 
$\Pr\left[\liminf_{t \to \infty} \frac{1}{t} (X_t - X_0) \geq - B^\mathrm{lower} \right] = 1$.
% \begin{align}
% \Pr\left[\liminf_{t \to \infty} \frac{1}{t} (X_t - X_0) \geq - B^\mathrm{lower} \right] = 1\label{eq:prop:lower}
% .
% \end{align}
\end{proposition}

\section{Proof}

\subsection{Proof of \Cref{prop:convergence_rate}}\label{apdx:prop:convergence_rate}
Let $Z_{t+1} = X_{t+1} - \E[X_{t+1} \mid \mathcal{F}_t]$.
Then, $\{Z_{t}\}_{t \geq 1}$ is a martingale difference sequence adapted to $\{\mathcal{F}_{t}\}$, and
\begin{multline}
\frac1t (X_t - X_0)
= \frac1t \sum_{i=1}^{t}(X_i - X_{i-1})
= \frac1t \sum_{i=1}^{t}(Z_i + \E[X_i - X_{i-1} \mid \mathcal{F}_{i-1}])
\\
= \frac1t \sum_{i=1}^{t}Z_i + \frac1t \sum_{i=1}^{t} \E[X_i - X_{i-1} \mid \mathcal{F}_{i-1}]
.
\end{multline}
Suppose that C3 holds.
{\color{blue}Beucause $\Var[X_t] = \Var[Z_t] + \Var[\E[X_t\mid\mathcal{F}_{t-1}]] \geq \Var[Z_t]$, }
C3 immediately implies that $\sum_{t=1}^{\infty}{\color{blue}\E[Z_{t}^2]} / t^2 < \infty$. Consequently, from the strong law of large numbers of martingale \cite{chow1967strong}, we obtain $\lim_{t\to\infty} \frac1t \sum_{i=1}^{t} Z_t = 0$ almost surely. 
We obtain statement (I) by taking $\limsup$ of the equation above and using C1 and $\limsup_{t\to\infty} \frac1t \sum_{i=1}^{t} Z_t = 0$.
Similarly, we obtain statement (II) by taking $\liminf$ and using C2 and $\liminf_{t\to\infty} \frac1t \sum_{i=1}^{t} Z_t = 0$. This completes the proof.

\begin{thebibliography}{99}
\bibitem[7]{chow1967strong} Yuan~Shih Chow et~al.
\newblock On a strong law of large numbers for martingales.
\newblock {\em The Annals of Mathematical Statistics}, 38(2):610--610, 1967.
\end{thebibliography}
\end{document}
